% GENERATED FILE. Edit canonical lessons, then run npm run generate:books.
% canonical-source-hash: fnv1a64:a7136f24b05fe9cd
% canonical-lessons: TE-C133-guddalu, TE-C133-misalu, TE-C133-kantiniru, TE-C133-padava, TE-C133-railu

\chapter{Clothes, Moustache, Tears, Boat, Train}
\label{ch:te-clothes-moustache-tears-boat-train}
\hlchaptermodality{133}

\begin{chapteropening}
\textbf{By the end of this chapter:} \emph{I can say clothes, a moustache, tears, a boat and a train.}

Complete the last lesson of chapter 133: i can say clothes, a moustache, tears, a boat and a train.
\end{chapteropening}

\section[\texorpdfstring{\te{గుడ్డలు} (guḍḍalu)}{గుడ్డలు (guḍḍalu)}]{\te{గుడ్డలు} (guḍḍalu) — clothes, cloth}

\label{lesson:TE-C133-guddalu}



\begin{quote}
Before the new one: say the Telugu for a blouse, then the Telugu for a lotus.
\end{quote}

\subsection*{You'll want to know: \te{గుడ్డలు}}
\textbf{\te{గుడ్డలు}} — \emph{guḍḍalu} — \textquotedblleft{}clothes, cloth\textquotedblright{}.

\begin{grammarlens}[title={What you've built}]
One more word for this chapter.
\end{grammarlens}

\subsection*{Your turn}
\begin{itemize}
\raggedright
  \item \emph{Say it:} \emph{guḍḍalu}
  \item \emph{Say it:} \emph{guḍḍalu}, once more
  \item \emph{Say it:} say \emph{guḍḍalu}
  \item \emph{Recall:} say the Telugu for a blouse, then the Telugu for a lotus, then say \emph{guḍḍalu} again
\end{itemize}

\begin{tcolorbox}[breakable,colback=teal!4,colframe=teal!35!black,title={Before you move on}]
What does \te{గుడ్డలు} mean? (\textquotedblleft{}Clothes, cloth\textquotedblright{}.) Say it once more.
\end{tcolorbox}
\section[\texorpdfstring{\te{మీసాలు} (mīsālu)}{మీసాలు (mīsālu)}]{\te{మీసాలు} (mīsālu) — a moustache}

\label{lesson:TE-C133-misalu}



\begin{quote}
Before the new one: say the Telugu for a lotus, then the Telugu for clothes.
\end{quote}

\subsection*{You'll want to know: \te{మీసాలు}}
\textbf{\te{మీసాలు}} — \emph{mīsālu} — \textquotedblleft{}a moustache\textquotedblright{}.

\begin{grammarlens}[title={What you've built}]
One more word for this chapter.
\end{grammarlens}

\subsection*{Your turn}
\begin{itemize}
\raggedright
  \item \emph{Say it:} \emph{mīsālu}
  \item \emph{Say it:} \emph{mīsālu}, once more
  \item \emph{Say it:} say \emph{mīsālu}
  \item \emph{Recall:} say the Telugu for a lotus, then the Telugu for clothes, then say \emph{mīsālu} again
\end{itemize}

\begin{tcolorbox}[breakable,colback=teal!4,colframe=teal!35!black,title={Before you move on}]
What does \te{మీసాలు} mean? (\textquotedblleft{}A moustache\textquotedblright{}.) Say it once more.
\end{tcolorbox}
\section[\texorpdfstring{\te{కంటినీరు} (kaṇṭinīru)}{కంటినీరు (kaṇṭinīru)}]{\te{కంటినీరు} (kaṇṭinīru) — tears}

\label{lesson:TE-C133-kantiniru}



\begin{quote}
Before the new one: say the Telugu for clothes, then the Telugu for a moustache.
\end{quote}

\subsection*{You'll want to know: \te{కంటినీరు}}
\textbf{\te{కంటినీరు}} — \emph{kaṇṭinīru} — \textquotedblleft{}tears\textquotedblright{}.

\begin{grammarlens}[title={What you've built}]
One more word for this chapter.
\end{grammarlens}

\subsection*{Your turn}
\begin{itemize}
\raggedright
  \item \emph{Say it:} \emph{kaṇṭinīru}
  \item \emph{Say it:} \emph{kaṇṭinīru}, once more
  \item \emph{Say it:} say \emph{kaṇṭinīru}
  \item \emph{Recall:} say the Telugu for clothes, then the Telugu for a moustache, then say \emph{kaṇṭinīru} again
\end{itemize}

\begin{tcolorbox}[breakable,colback=teal!4,colframe=teal!35!black,title={Before you move on}]
What does \te{కంటినీరు} mean? (\textquotedblleft{}Tears\textquotedblright{}.) Say it once more.
\end{tcolorbox}
\section[\texorpdfstring{\te{పడవ} (paḍava)}{పడవ (paḍava)}]{\te{పడవ} (paḍava) — a boat}

\label{lesson:TE-C133-padava}



\begin{quote}
Before the new one: say the Telugu for a moustache, then the Telugu for tears.
\end{quote}

\subsection*{You'll want to know: \te{పడవ}}
\textbf{\te{పడవ}} — \emph{paḍava} — \textquotedblleft{}a boat\textquotedblright{}.

\begin{grammarlens}[title={What you've built}]
One more word for this chapter.
\end{grammarlens}

\subsection*{Your turn}
\begin{itemize}
\raggedright
  \item \emph{Say it:} \emph{paḍava}
  \item \emph{Say it:} \emph{paḍava}, once more
  \item \emph{Say it:} say \emph{paḍava}
  \item \emph{Recall:} say the Telugu for a moustache, then the Telugu for tears, then say \emph{paḍava} again
\end{itemize}

\begin{tcolorbox}[breakable,colback=teal!4,colframe=teal!35!black,title={Before you move on}]
What does \te{పడవ} mean? (\textquotedblleft{}A boat\textquotedblright{}.) Say it once more.
\end{tcolorbox}
\section[\texorpdfstring{\te{రైలు} (railu)}{రైలు (railu)}]{\te{రైలు} (railu) — a train}

\label{lesson:TE-C133-railu}



\begin{quote}
Before the new one: say the Telugu for tears, then the Telugu for a boat.
\end{quote}

\subsection*{You'll want to know: \te{రైలు}}
\textbf{\te{రైలు}} — \emph{railu} — \textquotedblleft{}a train\textquotedblright{}.

\begin{grammarlens}[title={What you've built}]
One more word for this chapter.
\end{grammarlens}

\subsection*{Your turn}
\begin{itemize}
\raggedright
  \item \emph{Say it:} \emph{railu}
  \item \emph{Say it:} \emph{railu}, once more
  \item \emph{Say it:} say \emph{railu}
  \item \emph{Recall:} say the Telugu for tears, then the Telugu for a boat, then say \emph{railu} again
\end{itemize}

\begin{tcolorbox}[breakable,colback=teal!4,colframe=teal!35!black,title={Before you move on}]
What does \te{రైలు} mean? (\textquotedblleft{}A train\textquotedblright{}.) Say it once more.
\end{tcolorbox}
